\chapter[Hierarchy, Ordinals, and Cardinals]{The Cumulative Hierarchy, Ordinals, and Cardinals}
\label{ch:cumulative-hierarchy}

\section{Transitive sets and ordinals}
\label{sec:cumulative-hierarchy-01}

Transitive sets and ordinals were defined in
\cref{def:transitive-set,def:ordinal-cardinal}.  This chapter uses those exact
definitions inside set theory rather than introducing new terms for them.

For ordinals \(\alpha\) and \(\beta\), write \(\alpha<\beta\) iff
\(\alpha\in\beta\).  Define
\[
  \alpha+1\defeq\alpha\cup\set{\alpha}.
\]
An ordinal is a \emph{successor ordinal} iff it has this form; a nonzero ordinal
that is not a successor ordinal is a \emph{limit ordinal}.  These terms name
ordinal forms, not arithmetic operations on arbitrary sets.

\begin{theorem}[Trichotomy and transitivity of ordinals]
For ordinals \(\alpha\) and \(\beta\), exactly one of
\[
  \alpha\in\beta,\qquad \alpha=\beta,\qquad \beta\in\alpha
\]
holds.  Every member of an ordinal is an ordinal.
\end{theorem}

\paragraph{Proof sketch.}
The route runs through \(\alpha\cap\beta\).  First, \(\alpha\cap\beta\) is an
ordinal: it is transitive as an intersection of transitive sets, and \(\in\)
strictly well-orders it as a subset of \(\alpha\).  Second, an ordinal that
is a proper subset of an ordinal is a member of it: if
\(\gamma\subsetneq\delta\), the well-order of \(\delta\) supplies the
\(\in\)-least element \(\xi\) of \(\delta\setminus\gamma\); minimality of
\(\xi\) together with transitivity of \(\delta\) puts every member of
\(\xi\) in \(\gamma\), and linearity of \(\in\) on \(\delta\) together with
transitivity of \(\gamma\) puts every member of \(\gamma\) in \(\xi\), so
\(\gamma=\xi\in\delta\).  Now if \(\alpha\cap\beta\) were a proper subset of
both \(\alpha\) and \(\beta\), the second step would give
\(\alpha\cap\beta\in\alpha\cap\beta\), contradicting the irreflexivity of
the strict well-order on the ordinal \(\alpha\cap\beta\); hence
\(\alpha\cap\beta=\alpha\) or \(\alpha\cap\beta=\beta\), and the second step
converts the surviving inclusion, when proper, into \(\alpha\in\beta\) or
\(\beta\in\alpha\).  At most one alternative holds: any two of them combine,
using transitivity of the sets involved, into an ordinal that is a member of
itself.  For the final claim, a member \(x\) of an ordinal \(\alpha\) is a
subset of \(\alpha\) by transitivity of \(\alpha\), so \(\in\) strictly
well-orders \(x\); and \(x\) is transitive because transitivity of
\(\alpha\) puts the members of members of \(x\) inside \(\alpha\), where the
transitivity of the strict order \(\in\) on \(\alpha\) yields the required
memberships.  The omitted obligations are the two elementwise membership
calculations in the second step; they are supplied by \cite{Jech2003}.

\begin{theorem}[Transfinite induction]
Let \(\varphi(\alpha,\bar p)\) be a formula.  If
\[
  \forall\alpha\,
  \left[
    \bigl(\forall\beta<\alpha\,\varphi(\beta,\bar p)\bigr)
    \to\varphi(\alpha,\bar p)
  \right],
\]
then \(\varphi(\alpha,\bar p)\) holds for every ordinal \(\alpha\).
\end{theorem}

\begin{proof}
If the conclusion fails at some ordinal \(\delta\), separation applied to
\(\delta+1\) produces the nonempty set
\(\set{\beta\in\delta+1:\neg\varphi(\beta,\bar p)}\).  The ordinal
well-order supplies its least member \(\alpha\).  Every \(\beta<\alpha\) satisfies \(\varphi\), so
the displayed hypothesis gives \(\varphi(\alpha,\bar p)\), a contradiction.
\end{proof}

\begin{theorem}[Transfinite recursion]
Let \(G\) be a definable class operation assigning one set \(G(f)\) to every
function \(f\) whose domain is an ordinal.  There is a unique class function
\(F\) on the ordinals such that
\[
  F(\alpha)=G(F\restr\alpha)
\]
for every ordinal \(\alpha\).  For each ordinal \(\gamma\), the restriction
\(F\restr\gamma\) is a set.
\end{theorem}

\paragraph{Proof sketch.}
For each ordinal \(\gamma\), use transfinite induction to construct the unique
function \(f_\gamma\) with domain \(\gamma\) satisfying
\(f_\gamma(\alpha)=G(f_\gamma\restr\alpha)\).  At a successor stage append
the uniquely specified value; at a limit stage take the union of the coherent
earlier functions.  Replacement collects the earlier values and separation
forms the required graphs.  Transfinite induction proves
\(f_\beta=f_\gamma\restr\beta\) whenever \(\beta<\gamma\), so their union is
the asserted class function; the same induction proves uniqueness.  The
omitted Separation and Replacement instances are supplied in
\cite[Chapter 2]{Jech2003}.

Replacement is essential in the general theorem: it collects the values
already constructed at earlier stages.

